\begin{abstract}
Long-horizon world-model prediction can fail before rollout begins: a nominal
physical prior may be biased under parameter shift, whereas recent measurements
may be noisy or missing. We introduce the Complementary-Filtered Residual World
Model (\method), which applies a prespecified gain-0.5, mask-exact complementary
update over the observed context and initializes a probabilistic
physics-residual transition for 25-step open-loop forecasting. The contribution
is this auditable composition and its confirmatory evaluation, not a claim that
complementary or Kalman-style filtering is new or universally optimal. We train
\ConfirmatoryRuns{} models spanning five controlled nonlinear systems, three
clean/corrupted/parameter-shift regimes, eight independent seed-level
realizations, five broad comparators, and five gain/training-horizon
sensitivities. The frozen primary estimand is the mean of seed-specific relative
NRMSE-25 changes against a preselected physics-guided neural network style
comparator (PGNN-style),
with all 15 system--regime cells equally weighted. The estimated relative
improvement was \PrimaryImprovement\% (95\% bootstrap CI
[\PrimaryRelativeCILow\%, \PrimaryRelativeCIHigh\%]); the absolute
proposed-minus-PGNN-style difference was \PrimaryDifference{} (95\% CI
[\PrimaryCILow, \PrimaryCIHigh]), and the exact one-sided sign-flip value for a
5\% practical margin was $p=\PrimaryPValue$. Thus H1 was \PrimaryDecision,
but the prespecified H2 robustness claim was \RobustnessDecision{} because
clean noninferiority failed: corrupt/shifted cells improved by
\CorruptImprovement\%, whereas clean cells changed by \CleanImprovement\%.
The evidence is limited to transition prediction on a finite
synthetic-mechanistic benchmark with known future controls; it does not
establish policy quality, visual world modeling, or real-robot transfer.
\end{abstract}
